\documentclass[aspectratio=169,11pt]{beamer}
\usepackage[indonesian]{babel}
\usepackage{amsmath,amssymb,booktabs,graphicx,tikz,pgfplots,listings,xcolor}
\pgfplotsset{compat=1.18}
\usetikzlibrary{arrows.meta,positioning,calc,angles,quotes}
\usetheme{metropolis}
\setbeamertemplate{title page}{%
  {\usebeamerfont{title}\inserttitle\par}
  {\usebeamerfont{subtitle}\insertsubtitle\par}
  \usebeamertemplate{title separator}\medskip
  {\usebeamerfont{author}\insertauthor\par}
  {\usebeamerfont{date}\insertdate\par}}
\usefonttheme{professionalfonts}
\metroset{progressbar=frametitle,block=fill}
\setbeamerfont{title}{size=\Large}
\setbeamerfont{subtitle}{size=\normalsize}
\usebeamercolor{normal text}\colorlet{codeink}{normal text.fg}
\usebeamercolor{alerted text}\colorlet{codeaccent}{alerted text.fg!60!codeink}
\lstset{language={[LaTeX]TeX},basicstyle=\ttfamily\scriptsize\color{codeink},
  texcsstyle=*\bfseries\color{codeaccent},columns=fullflexible,keepspaces=true,
  showstringspaces=false,backgroundcolor=\color{codeink!6},frame=l,
  rulecolor=\color{codeaccent},xleftmargin=5pt,aboveskip=4pt,belowskip=4pt}
\lstdefinestyle{terminal}{language=bash,morekeywords={latexmk},keywordstyle=\bfseries\color{codeaccent}}
\newcommand{\tugas}[1]{\par\medskip\alert{\textbf{Coba sendiri}}\quad #1}
\title{Membuat presentasi akademik dengan Beamer}
\subtitle{Hari 2 --- Struktur, visualisasi, desain, dan praktik}
\author[Pelatihan UT dan UGM]{Pelatihan Universitas Terbuka\\[2pt]Departemen Matematika FMIPA UGM}
\date{Yogyakarta, 20 November 2026}

\begin{document}
\begin{frame}{Berapa lama satu getaran ayunan?}
Pada panjang tali 0,40 m, waktu sepuluh getaran adalah 12,7 s.
Berapa lama satu getaran?
\begin{enumerate}
\item<1-> Tentukan besaran yang dicari, yaitu periode $T$.
\item<2-> Bagi waktu total dengan jumlah getaran.
\item<3-> $T=12{,}7/10=1{,}27$ s.
\end{enumerate}
\small Majukan PDF satu halaman. Perhatikan judul dan nomor frame.
\end{frame}
  \begin{frame}[plain]\titlepage  \end{frame}
\section{Sesi 1. Struktur dasar Beamer}
\begin{frame}{Komponen sebuah frame}
\begin{columns}[T,onlytextwidth]
\begin{column}{.48\textwidth}
\textbf{Yang dibaca audiens}\\[3pt]
Judul frame menyatakan pertanyaan. Isi memuat teks, rumus, tabel, atau gambar. Latar belakang berada di balik keduanya.
\end{column}
\begin{column}{.48\textwidth}
\textbf{Yang membantu orientasi}\\[3pt]
Headline dan footline mengisi tepi atas dan bawah. Tema juga dapat menampilkan sidebar kiri dan kanan, bilah navigasi, simbol navigasi, serta logo.
\end{column}
\end{columns}
\bigskip
Pada contoh ayunan tadi, judul tetap di atas dan nomor frame tetap di bawah ketika langkah perhitungan muncul.
\par\medskip{\small\emph{Beamer User Guide}, \S8.2, hlm.~62.}
\end{frame}
\begin{frame}{Urutan Beamer menggambar komponen frame}
\begin{columns}[c,onlytextwidth]
\begin{column}{.53\textwidth}
\centering
\begin{tikzpicture}[x=1cm,y=1cm]
% Tiap bidang bergeser ke kanan. Sumbu x lokal miring ke atas.
\foreach \k in {1,...,13}{
  \pgfmathsetmacro{\offset}{.38*(\k-1)}
  \begin{scope}[shift={(\offset,0)},yslant=.54]
    \ifnum\k=2 \fill[orange!12] (0,0) rectangle (2.15,2.1);\fi
    \ifnum\k=3 \draw[step=.42,codeink!25,very thin] (0,0) grid (2.15,2.1);\fi
    \ifnum\k=4 \fill[teal!55] (0,0) rectangle (.37,2.1);\fi
    \ifnum\k=5 \foreach \h in {.65,.82,.99}{\draw[teal!80!black,thick] (.03,\h)--(.34,\h);}\fi
    \ifnum\k=6 \fill[teal!55] (1.78,0) rectangle (2.15,2.1);\fi
    \ifnum\k=7 \foreach \h in {1.35,1.52,1.69}{\draw[teal!80!black,thick] (1.81,\h)--(2.12,\h);}\fi
    \ifnum\k=8 \foreach \u in {1.45,1.64,1.83}{\fill[violet] (\u,.12) circle[radius=.035];}\fi
    \ifnum\k=9 \fill[blue!65] (.08,1.58) rectangle (.43,1.99);\fi
    \ifnum\k=10 \fill[violet!55] (0,1.94) rectangle (2.15,2.1);\fi
    \ifnum\k=11 \foreach \h/\w in {1.27/1.51,1.08/1.31,.89/1.63}{\draw[codeink,thick] (.48,\h)--(\w,\h);}\fi
    \ifnum\k=12 \fill[violet!55] (0,0) rectangle (2.15,.13);\fi
    \ifnum\k=1 \draw[dashed,codeink!45] (0,0) rectangle (2.15,2.1);
    \else\ifnum\k=13 \draw[dashed,codeink!45] (0,0) rectangle (2.15,2.1);
    \else \draw[codeink!60,thin] (0,0) rectangle (2.15,2.1);\fi\fi
  \end{scope}
  \draw[codeink!60] (\offset,0)--++(0,-.22)
    node[below,inner sep=1pt,font=\scriptsize]{\k};
}
\end{tikzpicture}
\end{column}
\begin{column}{.45\textwidth}
\small
1. Hook latar LaTeX\\
2. Kanvas latar\\
3. Gambar atau pola latar\\
4--5. Kanvas dan isi sidebar kiri\\
6--7. Kanvas dan isi sidebar kanan\\
8. Simbol navigasi\\
9. Logo\\
10. Headline\\
11. Isi frame\\
12. Footline\\
13. Hook depan LaTeX
\end{column}
\end{columns}
\medskip{\small Digambar ulang dari Gambar~1, \emph{Beamer User Guide}, \S8.2. Diagram asli oleh Andrew Stacey.}
\end{frame}
\begin{frame}{Target Hari 2}
  \begin{itemize}
    \item Mengubah laporan eksperimen menjadi presentasi akademik singkat.
    \item Menyusun persamaan, tabel, gambar, grafik, dan diagram yang terbaca.
    \item Menggunakan overlay secara terencana, bukan sebagai hiasan.
    \item Menyesuaikan tema, warna, font, dan identitas institusi.
  \end{itemize}
  \end{frame}
\begin{frame}{Peta belajar Hari 2}
  \begin{enumerate}
    \item Struktur dasar Beamer
    \item Materi matematis dan visual
    \item Desain presentasi
    \item Hands-on: presentasi eksperimen ayunan
  \end{enumerate}
  \end{frame}

\begin{frame}{Sesi 1 --- Struktur dasar Beamer}
\large 08.30--09.45 \quad 75 menit
\medskip
\begin{itemize}
  \item 20 menit: kelas, preamble, title page, dan frame
  \item 20 menit: section, subsection, tema, dan navigasi
  \item 20 menit: blocks, lists, commands, dan environments
  \item 15 menit: latihan membuat kerangka presentasi
\end{itemize}
\end{frame}
\begin{frame}[fragile]{Dokumen Beamer minimum}
\begin{lstlisting}
\documentclass[aspectratio=169,11pt]{beamer}
\usetheme{metropolis}
\title{Eksperimen Ayunan Sederhana}
\author{Nama Peserta}
\date{20 November 2026}
\begin{document}
\begin{frame}[plain]\titlepage  \end{frame} % akhir frame
\begin{frame}{Tujuan}
  \begin{itemize}
    \item Mengukur periode ayunan.
    \item Mengestimasi percepatan gravitasi.
  \end{itemize}
  \end{frame} % akhir frame
\end{document}
\end{lstlisting}
\tugas{Kompilasi contoh dan ganti judul serta nama penulis.}
\end{frame}
\begin{frame}[fragile]{Frame, section, dan subsection}
\begin{lstlisting}
\section{Latar belakang}
\subsection{Pertanyaan}
\begin{frame}{Pertanyaan penelitian}
  Apakah periode bertambah ketika panjang tali bertambah?
  \end{frame} % akhir frame
\end{lstlisting}
\begin{itemize}
  \item Satu \emph{frame} adalah satu unit tampilan, bukan selalu satu ide kecil.
  \item Section membantu struktur dan navigasi; subsection digunakan seperlunya.
  \item Awali setiap sesi dengan satu \lstinline|\section{...}|.
\end{itemize}
\end{frame}
\begin{frame}[fragile]{Blocks dan lists}
\begin{lstlisting}
\begin{block}{Model}
  \[T=2\pi\sqrt{L/g}\]
\end{block}
\begin{alertblock}{Perhatian}
  Model berlaku untuk sudut kecil.
\end{alertblock}
\begin{enumerate}
  \item Ukur panjang tali.
  \item Catat waktu sepuluh getaran.
\end{enumerate}
\end{lstlisting}
Gunakan block untuk mengelompokkan gagasan. Batasi jumlah tingkat daftar agar slide tetap ringan.
\end{frame}
\begin{frame}{Tema dan navigasi}
\begin{itemize}
  \item Tema mengatur warna, font, block, dan struktur visual.
  \item Gunakan satu tema yang konsisten untuk seluruh pelatihan.
  \item Navigasi otomatis berasal dari section dan frame.
  \item Jika simbol navigasi tidak membantu, sembunyikan secara konsisten.
\end{itemize}
\begin{block}{Prinsip kerja}
  Preamble bersama membuat identitas semua presentasi tetap seragam.
\end{block}
\end{frame}

\begin{frame}[fragile]{Hubungkan bagian presentasi dengan daftar isi}
Letakkan daftar isi sesudah halaman judul. Tambahkan bagian dan subbagian
sebelum frame yang bersangkutan, lalu kompilasi ulang.
\begin{lstlisting}
\begin{frame}{Alur eksperimen}
  \tableofcontents
\end{frame} % akhir frame
\section{Metode}
\subsection{Pengukuran waktu}
\end{lstlisting}
Metropolis menampilkan pemisah saat bagian baru dimulai. Di pembaca PDF,
klik entri pada daftar isi untuk melompat ke bagian itu.
\end{frame}
\begin{frame}[fragile]{Simpan pengaturan yang dipakai setiap kelompok}
Simpan tiga definisi berikut dalam \texttt{gaya.tex}. Panggil
\lstinline|\input{gaya}| di prambule setelah memilih tema.
\begin{lstlisting}
\setbeamertemplate{navigation symbols}{}
\newcommand{\periode}{\ensuremath{T}}
\newenvironment{hasil}{\begin{exampleblock}{Hasil}}
                      {\end{exampleblock}}
\end{lstlisting}
Di dalam frame, gunakan perintah untuk simbol dan lingkungan untuk
mengelompokkan kesimpulan.
\begin{lstlisting}
\begin{hasil}
  Periode \periode{} meningkat ketika tali diperpanjang.
\end{hasil}
\end{lstlisting}
\end{frame}
\begin{frame}{Latihan pertama menghasilkan kerangka presentasi}
Buat proyek dengan satu berkas utama. Di Overleaf pilih XeLaTeX sebagai
mesin kompilasi. Secara lokal gunakan XeLaTeX melalui \texttt{latexmk}.
\begin{enumerate}
\item Beri judul berupa pertanyaan dari laporan Hari 1.
\item Buat dua bagian, yaitu metode dan hasil, masing-masing dengan satu frame.
\item Tambahkan daftar isi dan buka tautan menuju hasil.
\end{enumerate}
\tugas{Tukar PDF dengan rekan. Minta ia menemukan pertanyaan penelitian dan frame hasil.}
\end{frame}
\section{Sesi 2. Materi matematis dan visual}
\begin{frame}{Sesi 2 --- Materi matematis dan visual}
\large 10.00--11.30 \quad 90 menit
\medskip
\begin{itemize}
  \item 20 menit menulis persamaan dan meringkas tabel
  \item 15 menit menjelaskan cara ukur dengan gambar atau skema
  \item 25 menit memilih sumbu dan membuat grafik dari data
  \item 15 menit mengatur urutan pertanyaan dan jawaban dengan overlay
  \item 15 menit menyusun dua frame hasil dan mengujinya berpasangan
\end{itemize}
\end{frame}
\begin{frame}{Persamaan pada slide}
\begin{block}{Model periode ayunan}
  \[
    T\approx2\pi\sqrt{\frac{L}{g}},\qquad
    T^2\approx\frac{4\pi^2}{g}L.
  \]
\end{block}
\begin{itemize}
  \item Tampilkan hanya persamaan yang diperlukan untuk argumen.
  \item Definisikan simbol dan satuan di dekat persamaan.
  \item Hindari menyalin seluruh derivasi laporan ke satu frame.
\end{itemize}
\end{frame}
\begin{frame}[fragile]{Bagaimana meratakan tanda sama dengan?}
Gunakan \lstinline|aligned| dalam mode matematika. Tanda
\lstinline|&| menempatkan ketiga tanda sama dengan pada satu garis vertikal.
\begin{columns}[c,onlytextwidth]
  \begin{column}{.56\textwidth}
\begin{lstlisting}
\[
\begin{aligned}
  T &= t/N \\
    &= 15{,}6\,\mathrm{s}/10 \\
    &= 1{,}56\,\mathrm{s}.
\end{aligned}
\]
\end{lstlisting}
  \end{column}
  \begin{column}{.40\textwidth}
    \[
    \begin{aligned}
      T &= t/N \\
        &= 15{,}6\,\mathrm{s}/10 \\
        &= 1{,}56\,\mathrm{s}.
    \end{aligned}
    \]
  \end{column}
\end{columns}
\tugas{Ganti waktu sepuluh getaran menjadi 18,0 s. Perbarui hasil dan cocokkan dengan tabel pada halaman berikut.}
\end{frame}
\begin{frame}{Tabel: tampilkan pesan datanya}
\centering
\small
\begin{tabular}{rrrr}
\toprule
Panjang $L$ (m) & Getaran & Waktu (s) & Periode $T$ (s)\\
\midrule
0,40 & 10 & 12,7 & 1,27\\
0,60 & 10 & 15,6 & 1,56\\
0,80 & 10 & 18,0 & 1,80\\
\bottomrule
\end{tabular}
\medskip

Data ilustratif Hari 1. Untuk membandingkan periode, cukup tampilkan $L$ dan $T$.
\end{frame}
\begin{frame}[fragile]{Kolom mana yang menjawab pertanyaan audiens?}
Berapa periode saat panjang tali 0,60 m? Muat \lstinline|\usepackage{booktabs}|
di prambule, lalu tempatkan tabel ini di dalam frame.
\begin{columns}[c,onlytextwidth]
  \begin{column}{.55\textwidth}
\begin{lstlisting}
\begin{tabular}{rr}
\toprule
$L$ (m) & $T$ (s) \\
\midrule
0,40 & 1,27 \\
0,60 & \textbf{1,56} \\
0,80 & 1,80 \\
\bottomrule
\end{tabular}
\end{lstlisting}
  \end{column}
  \begin{column}{.41\textwidth}
    \centering
    \begin{tabular}{rr}
      \toprule $L$ (m) & $T$ (s)\\
      \midrule 0,40 & 1,27\\ 0,60 & \textbf{1,56}\\ 0,80 & 1,80\\
      \bottomrule
    \end{tabular}
    \par\smallskip{\small Data ilustratif Hari 1.}
  \end{column}
\end{columns}
Satuan di kepala kolom berlaku untuk semua baris. Tebalkan nilai yang akan disebut saat presentasi.
\end{frame}
\begin{frame}{Diagram dengan TikZ}
\centering
\begin{tikzpicture}[scale=1.15,>=Latex]
  \coordinate (O) at (0,2);
  \coordinate (B) at ($(O)+(-65:2.5)$);
  \draw[thick] (-1,2)--(1,2);
  \draw[dashed] (O)--++(0,-2.8);
  \draw[thick] (O)--(B) node[midway,right] {$L$};
  \draw (O)++(-90:.75) arc (-90:-65:.75);
  \node at ($(O)+(-77:1.05)$) {$\theta$};
  \fill[alerted text.fg] (B) circle (.14);
\end{tikzpicture}
\qquad
\begin{minipage}{.34\linewidth}
  \small Diagram harus memperjelas relasi, bukan sekadar menghias slide.\end{minipage}
\end{frame}
\begin{frame}[fragile]{Tampilkan foto alat pada lebar kolomnya}
Simpan foto peserta sebagai \texttt{gambar/alat.jpg}. Tempatkan kode ini
di dalam frame metode.
\begin{lstlisting}
\begin{columns}[c,onlytextwidth]
  \begin{column}{.48\textwidth}
    \includegraphics[width=\linewidth,height=4cm,
                     keepaspectratio]{gambar/alat.jpg}
  \end{column}
  \begin{column}{.48\textwidth}
    Ukur dari titik gantung ke pusat beban.
  \end{column}
\end{columns}
\end{lstlisting}
\tugas{Tunjukkan kedua ujung pengukuran pada foto. Jika alat sulit terlihat, gunakan skema pada contoh praktik.}
\end{frame}
\begin{frame}[fragile]{Bangun skema dari titik gantung, tali, dan beban}
\begin{lstlisting}
\begin{tikzpicture}
  \draw[thick] (-1,0)--(1,0);
  \draw[dashed] (0,0)--(0,-2.7);
  \draw[thick] (0,0)--(-72:2.6) node[midway,right] {$L$};
  \fill (-72:2.6) circle (.14);
\end{tikzpicture}
\end{lstlisting}
Garis putus-putus menandai posisi setimbang. Tali dan pusat beban
bertemu pada koordinat yang sama.
\tugas{Ubah sudut pada dua koordinat polar dan periksa titik akhir tali.}
\end{frame}
\begin{frame}[fragile]{Agar label mengikuti beban ketika sudut diubah}
Simpan titik gantung dan pusat beban sebagai koordinat bernama. Ujung tali,
lingkaran, serta panah label kemudian memakai titik yang sama.
\begin{columns}[c,onlytextwidth]
\begin{column}{.57\textwidth}
\begin{lstlisting}
\begin{tikzpicture}
  \coordinate (O) at (0,0);
  \coordinate (B) at (-70:2);
  \draw[thick] (O)--(B);
  \fill (B) circle (.12);
  \draw[<-] (B)++(.13,0)--++(.55,0)
    node[right]{pusat beban};
\end{tikzpicture}
\end{lstlisting}
\end{column}
\begin{column}{.40\textwidth}
\centering
\begin{tikzpicture}
  \coordinate (O) at (0,0);
  \coordinate (B) at (-70:2);
  \draw[thick] (O)--(B);
  \fill (B) circle (.12);
  \draw[<-] (B)++(.13,0)--++(.55,0) node[right]{pusat beban};
\end{tikzpicture}
\end{column}
\end{columns}
\tugas{Ganti sudut menjadi $-50^\circ$ hanya pada definisi B. Periksa kembali letak panah.}
\end{frame}
\begin{frame}[fragile]{Pilihan diagram sains}
\begin{itemize}
  \item TikZ: diagram konseptual dan geometri.
  \item CircuiTikZ: rangkaian listrik dan simbol komponen.
  \item Feynman diagrams: proses interaksi partikel.
  \item Paket khusus dimuat hanya ketika memang diperlukan.
\end{itemize}
\begin{lstlisting}
% Dengan \usepackage{circuitikz} di prambule:
\begin{circuitikz}
  \draw (0,0) to[R,l=$R$] (2,0) to[C,l=$C$] (4,0);
\end{circuitikz}
\end{lstlisting}
Pada rangkaian RC seri ini, label menunjukkan komponen yang dilalui arus.
\end{frame}
\begin{frame}[fragile]{Pilihan fisika partikel untuk kelas yang memerlukannya}
Pada diagram hamburan elektron, garis berpanah menunjukkan fermion.
Garis foton menghubungkan dua titik interaksi. Contoh berikut memakai
paket \lstinline|tikz-feynman| dengan posisi titik ditentukan manual.
\begin{lstlisting}
\begin{tikzpicture}
  \begin{feynman}
    \vertex (a) at (0,1); \vertex (b) at (0,-1);
    \vertex (i) at (-1,1.6);  \vertex (j) at (1,1.6);
    \vertex (k) at (-1,-1.6); \vertex (l) at (1,-1.6);
    \diagram* {
      (i) -- [fermion] (a) -- [fermion] (j),
      (k) -- [fermion] (b) -- [fermion] (l),
      (a) -- [photon] (b)
    };
  \end{feynman}
\end{tikzpicture}
\end{lstlisting}
\end{frame}
\begin{frame}{Satu tabel dapat menjawab dua pertanyaan grafik}
\begin{columns}[T,onlytextwidth]
\begin{column}{.49\textwidth}
\centering\textbf{Berapa periode pada setiap panjang?}\par
\begin{tikzpicture}
\begin{axis}[width=\linewidth,height=4.1cm,xlabel={$L$ (m)},ylabel={$T$ (s)},
  xmin=.32,xmax=.88,xtick={.4,.6,.8},ymin=1.1,ymax=2]
  \addplot[only marks,mark=*] coordinates {(.4,1.27) (.6,1.56) (.8,1.80)};
\end{axis}
\end{tikzpicture}
\end{column}
\begin{column}{.49\textwidth}
\centering\textbf{Bagaimana menguji model ayunan?}\par
\begin{tikzpicture}
\begin{axis}[width=\linewidth,height=4.1cm,xlabel={$L$ (m)},ylabel={$T^2$ (s$^2$)},
  xmin=.32,xmax=.88,xtick={.4,.6,.8},ymin=1.35,ymax=3.6]
  \addplot[domain=.37:.83,no marks] {4.05*x};
  \addplot[only marks,mark=*] coordinates {(.4,1.613) (.6,2.434) (.8,3.240)};
\end{axis}
\end{tikzpicture}
\end{column}
\end{columns}
Di kiri, baca periode langsung dalam sekon. Di kanan, bandingkan titik dengan garis acuan $T^2=4{,}05L$ yang mengikuti bentuk model sudut kecil.
\end{frame}
\begin{frame}{Visualisasi data dengan pgfplots}
\centering
\begin{tikzpicture}
\begin{axis}[width=.72\linewidth,height=5.2cm,grid=major,
  xlabel={$L$ (m)},ylabel={$T^2$ (s$^2$)},legend pos=north west]
\addplot[only marks,mark=*] coordinates {(0.4,1.613) (0.6,2.434) (0.8,3.240)};
\addplot[no markers,domain=.35:.85] {4.05*x};
\legend{Data ilustratif,Garis $T^2=4{,}05L$}
\end{axis}
\end{tikzpicture}
\end{frame}
\begin{frame}[fragile]{Plot pasangan panjang dan kuadrat periode}
Muat \lstinline|\usepackage{pgfplots}| dan tetapkan
\lstinline|\pgfplotsset{compat=1.18}| di prambule.
\begin{lstlisting}
\begin{tikzpicture}
  \begin{axis}[height=5cm, xlabel={$L$ (m)},
               ylabel={$T^2$ (s$^2$)}]
    \addplot[only marks] coordinates
      {(0.4,1.613) (0.6,2.434) (0.8,3.240)};
    \addplot[domain=.35:.85] {4.05*x};
  \end{axis}
\end{tikzpicture}
\end{lstlisting}
Titik pertama berasal dari $1{,}27^2\approx1{,}613$. Gunakan titik
desimal di dalam koordinat sumber.
\end{frame}
\begin{frame}[fragile]{Hitung kuadrat periode langsung dari tabel sumber}
Di dalam lingkungan \lstinline|axis| untuk grafik $T^2$ terhadap $L$,
ganti perintah yang menggambar titik dengan potongan berikut.
\begin{lstlisting}
\addplot[only marks,mark=*]
  table[x=L, y expr={\thisrow{T}^2}] {
    L     T
    0.40  1.27
    0.60  1.56
    0.80  1.80
  };
\end{lstlisting}
Untuk panjang 0,60 m, ekspresi pada sumbu vertikal menghasilkan
$1{,}56^2=2{,}4336$ s$^2$.
\tugas{Ubah periode baris tengah menjadi 1,60 s pada salinan contoh.
Prediksi titik yang bergeser dan nilai barunya sebelum kompilasi.}
\end{frame}
\begin{frame}{Apa yang dapat kita baca dari kemiringan garis?}
Dua titik terluar memberi perkiraan gradien
\[
 m=\frac{3{,}240-1{,}613}{0{,}80-0{,}40}
 \approx4{,}07\ \mathrm{s^2/m}.
\]
Garis acuan pada grafik memakai $m=4{,}05$ s$^2$/m. Dari model
$T^2\approx(4\pi^2/g)L$, garis itu memberi $g\approx9{,}75$ m/s$^2$.
\tugas{Gunakan tabel untuk menemukan satu periode. Gunakan grafik untuk menjelaskan hubungan ketiga pengukuran.}
\end{frame}
\begin{frame}[fragile]{Overlay dan incremental reveal}
\begin{lstlisting}
\begin{frame}{Dari pengukuran ke kesimpulan}
  \begin{enumerate}
    \item<1-> Sepuluh getaran memerlukan 12,7 s.
    \item<2-> Periode satu getaran adalah 1,27 s.
    \item<3-> Nilai tersebut dibandingkan dengan model.
  \end{enumerate}
\end{frame} % akhir frame
\end{lstlisting}
\begin{itemize}
  \item Gunakan overlay untuk mengatur urutan penjelasan.
  \item Pastikan slide tetap masuk akal pada setiap tahap.
\end{itemize}
\end{frame}
\begin{frame}[fragile]{Pertahankan posisi pertanyaan saat jawaban muncul}
\begin{lstlisting}
Berapa periode untuk sepuluh getaran dalam 12,7 s?
\uncover<2->{Jawabannya adalah $1{,}27$ s.}
\end{lstlisting}
\lstinline|\uncover<2->| menyisakan ruang sejak halaman pertama.
\lstinline|\only<2->| menyisipkan isi mulai halaman kedua sehingga
posisi teks lain dapat bergeser.
\tugas{Tambahkan satu baris teks di bawah jawaban. Coba kedua perintah dan bandingkan posisi baris itu.}
\end{frame}
\begin{frame}{Prediksi dahulu, lalu buka hasil pengukuran}
\begin{columns}[c,onlytextwidth]
  \begin{column}{.50\textwidth}
    \textbf{Jika panjang tali dilipatduakan, apakah periode juga dua kali lipat?}
    \par\medskip Pada panjang 0,40 m, periodenya 1,27 s.
    Berikan prediksi untuk panjang 0,80 m.
  \end{column}
  \begin{column}{.46\textwidth}
    \uncover<2->{\begin{block}{Data ilustratif}
      Pada 0,80 m, periodenya 1,80 s. Rasio kedua periode sekitar
      $1{,}80/1{,}27\approx1{,}42$.
    \end{block}}
    \uncover<3->{Model sudut kecil memprediksi faktor $\sqrt{2}\approx1{,}41$.}
  \end{column}
\end{columns}
\end{frame}
\begin{frame}[fragile]{Atur kapan audiens melihat bukti dan penjelasan}
\begin{lstlisting}
Jika panjang tali dua kali lipat, bagaimana periodenya?
\uncover<2->{Data memberi rasio $1{,}80/1{,}27$.}
\uncover<3->{Model memprediksi faktor $\sqrt{2}$.}
\end{lstlisting}
Spesifikasi \lstinline|<2->| menampilkan bukti mulai langkah kedua
dan mempertahankannya saat penjelasan muncul pada langkah ketiga.
\tugas{Ubah spesifikasi bukti menjadi \lstinline|<2>|. Majukan PDF sampai
langkah ketiga. Pilih versi yang memungkinkan audiens membandingkan data dengan model.}
\end{frame}
\begin{frame}{Praktik 15 menit untuk menyajikan satu hasil eksperimen}
\begin{enumerate}
\item \textbf{5 menit.} Buat frame pengukuran untuk $L=0{,}60$ m.
      Tampilkan waktu sepuluh getaran. Beri audiens kesempatan menghitung
      $T$ sebelum membuka jawabannya dengan overlay.
\item \textbf{7 menit.} Buat frame hasil dengan tiga titik $T^2$ terhadap $L$.
      Beri kedua sumbu satuan. Tulis judul yang menyatakan kesimpulan dari ketiga titik.
\item \textbf{3 menit.} Tukarkan PDF. Minta rekan membaca nilai periode
      dan menunjuk titik yang sesuai pada grafik. Revisi label yang membuatnya ragu.
\end{enumerate}
\small\textbf{Cek hasil}\quad Pada 0,60 m, jawaban lisan adalah 1,56 s
dan ordinat titik sekitar 2,43 s$^2$.
\end{frame}

\section{Sesi 3. Desain presentasi}
\begin{frame}{Sesi 3 --- Desain presentasi}
\large 13.00--14.15 \quad 75 menit
\medskip
\begin{itemize}
  \item Tipografi, ukuran teks, dan kontras
  \item Warna, tata letak, dan konsistensi visual
  \item Tema institusi dan preamble bersama
  \item Font modern, XeLaTeX, dan teks multilingual
\end{itemize}
\end{frame}
\begin{frame}[fragile]{Warna, font, dan identitas}
\begin{lstlisting}
\definecolor{laut}{HTML}{145A86}
\setbeamercolor{frametitle}{fg=white,bg=laut}
\setbeamerfont{frametitle}{series=\bfseries}
\setbeamertemplate{navigation symbols}{}
\end{lstlisting}
\begin{itemize}
  \item Pilih warna dengan kontras yang cukup.
  \item Jangan menggunakan warna sebagai satu-satunya penanda informasi.
  \item Dengan XeLaTeX, gunakan \lstinline|fontspec| untuk font modern.
\end{itemize}
\end{frame}
\begin{frame}{Prinsip slide matematika dan sains}
\begin{enumerate}
  \item Satu frame menyampaikan satu gagasan utama.
  \item Teks besar, kontras tinggi, dan ruang kosong cukup.
  \item Satuan, label sumbu, legenda, dan caption harus terbaca.
  \item Grafik dipilih untuk menjawab pertanyaan, bukan sekadar menampilkan data.
\end{enumerate}
\tugas{Periksa slide Anda dari jarak paling belakang ruangan.}
\end{frame}

\begin{frame}[fragile]{Pilih font dan uji karakter yang dipakai}
Di prambule setelah tema, pilih font sans serif yang tersedia pada
komputer penyusun. Kompilasi dengan XeLaTeX.
\begin{lstlisting}
\usepackage{fontspec}
\setsansfont{Latin Modern Sans}
\setbeamerfont{frametitle}{size=\large,series=\bfseries}
\end{lstlisting}
Uji judul Indonesia dan Inggris, nama Miletić, serta rumus $T^2$.
Huruf teks mengikuti font sans serif, sedangkan rumus memakai font matematika.
\tugas{Periksa aksen, pangkat, dan tanda minus pada PDF. Cari pesan \texttt{Missing character} pada log jika ada karakter yang hilang.}
\end{frame}
\begin{frame}{Uji desain dengan satu pertanyaan tentang data}
Proyeksikan tabel empat kolom dari Sesi 2. Minta rekan di belakang
ruangan menyebutkan periode pada $L=0{,}60$ m.
\begin{enumerate}
\item Sisakan kolom panjang dan periode. Pindahkan rincian waktu ke frame metode.
\item Tulis judul yang menyatakan pola. Cantumkan satuan di kepala kolom.
\item Gunakan tebal atau label untuk nilai yang dibahas. Periksa juga dalam skala abu-abu.
\end{enumerate}
Ulangi pertanyaan yang sama. Jawaban yang diharapkan adalah 1,56 s.
\end{frame}
\subsection{Frame standout pada tema Metropolis}
\begin{frame}[fragile]{Soroti satu temuan dengan frame standout}
Metropolis menyediakan opsi \lstinline|standout| untuk menempatkan
satu pesan pendek di tengah halaman dengan warna kontras. Frame ini
tidak menambah penghitung frame. Contoh berikut merangkum hasil ayunan.
\begin{lstlisting}
\begin{frame}[standout]
  Panjang tali dua kali lipat.\\
  Periode sekitar $1{,}42$ kali semula.
\end{frame} % akhir frame
\end{lstlisting}
\tugas{Pilih satu angka dari hasil kelompok. Tulis satu kalimat yang
menjelaskan maknanya, lalu tampilkan sesudah grafik pendukung.}
\end{frame}
\begin{frame}[standout]
  Panjang tali dua kali lipat.\\
  Periode sekitar $1{,}42$ kali semula.
\end{frame}

\section{Sesi 4. Hands-on Beamer}
\begin{frame}{Sesi 4 --- Hands-on Beamer}
\large 14.30--16.00 \quad 90 menit
\medskip
\begin{enumerate}
  \item Buat presentasi 5--7 menit dari laporan Hari 1.
  \item Masukkan persamaan, tabel, diagram, dan grafik.
  \item Tambahkan overlay pada satu bagian yang memerlukan urutan.
  \item Sesuaikan tema dan warna secara konsisten.
  \item Kompilasi, periksa PDF, lalu uji tampilan akhir.
\end{enumerate}
\end{frame}
\begin{frame}{Rancangan presentasi mini}
\begin{enumerate}
  \item Masalah dan tujuan
  \item Model matematis
  \item Metode dan diagram alat
  \item Data dan grafik $T^2$ terhadap $L$
  \item Estimasi $g$ dan keterbatasan
  \item Kesimpulan
\end{enumerate}
\end{frame}
\begin{frame}[fragile]{Pemeriksaan kompilasi dan tampilan}
\begin{lstlisting}[style=terminal]
latexmk -xelatex ayunan.tex
\end{lstlisting}
\begin{itemize}
  \item Tidak ada error yang tersisa.
  \item Tidak ada teks atau grafik yang keluar dari batas frame.
  \item Semua label sumbu, angka, dan simbol terbaca.
  \item Setiap frame mendukung narasi lisan.
\end{itemize}
\end{frame}
\begin{frame}{Alokasi praktik dan telaah berpasangan}
\begin{description}
\item[15 menit] Pilih pertanyaan, susun judul dan urutan 5--7 frame.
\item[35 menit] Masukkan rumus, diagram atau foto, tabel dan grafik.
  Gunakan data ilustratif dalam \texttt{ayunan.tex} atau data sendiri yang diberi sumber.
\item[15 menit] Atur satu overlay dan terapkan warna institusi di prambule.
\item[15 menit] Berpasangan, bergantian presentasi 5--7 menit. Pendengar menyebutkan kesimpulan dan angka pendukungnya.
\item[10 menit] Revisi bagian yang sulit dibaca. Serahkan PDF dan sumber beserta asetnya.
\end{description}
\end{frame}
\begin{frame}[fragile]{Temukan penyebab ketika hasil kompilasi berbeda}
\begin{description}
\item[Gambar hilang] Cocokkan jalur dan huruf besar-kecil nama berkas dengan argumen \lstinline|\includegraphics|.
\item[Listing kode] Gunakan \lstinline|\begin{frame}[fragile]| untuk frame berisi verbatim atau listing. Tulis penutup frame pada baris tersendiri.
\item[Daftar isi kosong] Kompilasi ulang agar daftar isi membaca informasi bagian terbaru.
\item[Isi terpotong] Cari \texttt{Overfull} pada log, lalu periksa frame terkait. Bagi tabel atau pindahkan rincian ke frame lain.
\end{description}
\end{frame}
\begin{frame}{Checklist akhir}
\begin{itemize}
  \item Struktur: title page, section, frame, dan navigasi.
  \item Isi: matematika, tabel, gambar, TikZ, pgfplots, overlay.
  \item Desain: font, ukuran, warna, tata letak, konsistensi.
  \item Praktik: kompilasi bersih dan pemeriksaan PDF.
\end{itemize}
\end{frame}
% \begin{frame}{Bacaan untuk mengembangkan contoh}
% \emph{The beamer class, User Guide}, versi 3.78, tersedia di folder
% \texttt{Resources/Beamer} pelatihan.
% \begin{itemize}
% \item Bab 3 dan 8 membahas presentasi pertama dan komponen frame.
% \item Bab 9 membahas overlay. Bab 10 membahas bagian dan daftar isi.
% \item Bagian 12.3, 12.6--12.7, dan 12.9 membahas blok, tabel, gambar, kolom, dan kode verbatim.
% \item Bab 5 serta bab 15--18 membahas keterbacaan, tema, templat, warna, dan font.
% \end{itemize}
% Untuk pengayaan, bandingkan kebutuhan presentasi dengan Powerdot
% dan kebutuhan satu bidang baca besar pada poster ilmiah.
% \end{frame}
\end{document}
